\chapter{Обзор и анализ методов безопасного обучения с подкреплением для автономных интеллектуальных агентов}
\label{ch:related_work}

Безопасное обучение с подкреплением направлено на построение автономных агентов, которые должны не только достигать целевых состояний и максимизировать награду, но и соблюдать заданные ограничения безопасности. В отличие от классической постановки обучения с подкреплением, где качество поведения определяется главным образом суммарной наградой, в безопасной постановке необходимо одновременно учитывать эффективность решения задачи и стоимость нарушений. Такие нарушения могут соответствовать столкновениям с препятствиями, попаданию в опасные области, превышению допустимого риска, повреждению объектов среды или невыполнению правил, заданных пользователем.

Ключевая трудность безопасного обучения состоит в противоречии между исследованием среды и соблюдением ограничений. Для обнаружения успешных траекторий агент должен пробовать разные действия и посещать новые состояния. Однако именно исследование часто приводит к нарушениям безопасности, особенно на ранних этапах обучения, когда стратегия ещё не обладает достоверной оценкой последствий своих действий. Эта проблема усиливается в задачах с редкой наградой, большим горизонтом планирования, частичной наблюдаемостью и сложной динамикой среды.

Существующие методы безопасного обучения решают эту проблему разными способами. Методы оптимизации стратегии с ограничениями вводят стоимость нарушений в задачу оптимизации и позволяют обучать стратегию без явного знания динамики среды. Однако они могут страдать от нестабильности двойственной оптимизации, консервативного поведения и слабого исследования в задачах с большим горизонтом. Методы безопасного управления и фильтрации действий усиливают контроль поведения на этапе исполнения, но часто требуют полной информации о состоянии среды, дифференцируемых ограничений, формальных правил или дополнительной оптимизации на каждом шаге.

Другим направлением являются методы на основе моделей мира. Они позволяют заранее оценивать последствия действий агента, строить воображаемые траектории и уменьшать число небезопасных взаимодействий со средой. Однако такие методы сталкиваются с ошибками предсказания, ограниченным горизонтом планирования и необходимостью корректно оценивать не только мгновенную стоимость, но и долгосрочный риск будущих нарушений. Поэтому для безопасного модельного обучения существенным становится вопрос построения функций долгосрочной стоимости и устойчивой агрегации их оценок.

Для задач с большим горизонтом важную роль играет иерархическое обучение с подкреплением. Иерархическая структура позволяет разбивать исходную задачу на более короткие подзадачи, использовать подцели или навыки и тем самым улучшать исследование среды. Однако классические иерархические методы не гарантируют безопасность: достижимая подцель может требовать небезопасной траектории. Поэтому в безопасном иерархическом обучении необходимо учитывать не только достижимость подцелей, но и стоимость их достижения, а также избегать чрезмерной зависимости от полной наблюдаемости, ручных правил и дополнительной фильтрации действий при исполнении.

Отдельная проблема возникает при переходе от фиксированных численных функций стоимости к естественно-языковым ограничениям. В практических сценариях пользователь часто задаёт цель и правила безопасности не в виде математической функции, а в форме текстовой инструкции. Это требует от агента связывать язык с объектами, действиями, историей взаимодействия и состояниями среды. Поэтому безопасное обучение с естественно-языковыми целями и ограничениями расширяет классическую постановку Safe RL и требует отдельных бенчмарков для оценки таких агентов.

В данной главе рассматриваются основные постановки и методы безопасного обучения с подкреплением. Сначала вводится формальная постановка через марковские процессы принятия решений с ограничениями. Далее анализируются методы оптимизации стратегии с ограничениями, методы безопасного управления и фильтрации действий, методы на основе моделей мира, иерархические методы для задач с большим горизонтом и безопасное иерархическое обучение. В конце главы рассматриваются бенчмарки безопасного обучения и переход к мультимодальным агентам с естественно-языковыми целями и ограничениями. Проведённый обзор позволяет выделить основные нерешённые проблемы, которые рассматриваются в последующих главах диссертации.

\section{Теоретические основы безопасного обучения с подкреплением}
\label{sec:safe_rl_formalization}

Классическое обучение с подкреплением рассматривает взаимодействие агента со средой как марковский процесс принятия решений:
\begin{equation}
\mathcal{M} = \langle \mathcal{S}, \mathcal{A}, P, R, \rho_0, \gamma \rangle,
\end{equation}
где $\mathcal{S}$ обозначает пространство состояний, $\mathcal{A}$~--- пространство действий, $P(s' \mid s,a)$~--- функцию переходов, $R(s,a)$~--- функцию награды, $\rho_0$~--- распределение начальных состояний, а $\gamma \in [0,1)$~--- коэффициент дисконтирования. Агент задаётся стратегией $\pi(a \mid s)$ и стремится максимизировать ожидаемую дисконтированную награду:
\begin{equation}
J_R(\pi) =
\mathbb{E}_{\tau \sim \pi}
\left[
\sum_{t=0}^{T} \gamma^t R(s_t,a_t)
\right].
\end{equation}

В задачах безопасного обучения с подкреплением дополнительно вводится функция стоимости или штрафа $C(s,a)$, описывающая нарушение ограничений. Такая постановка обычно формализуется как марковский процесс принятия решений с ограничениями:
\begin{equation}
\mathcal{M}_c =
\langle \mathcal{S}, \mathcal{A}, P, R, C, \rho_0, \gamma, \gamma_c, d \rangle,
\end{equation}
где $C$ задаёт стоимость небезопасного поведения, $\gamma_c$~--- коэффициент дисконтирования для стоимости, а $d$~--- допустимый бюджет нарушений. Целью агента является решение задачи:
\begin{equation}
\pi^* =
\arg\max_{\pi} J_R(\pi),
\quad
\text{при условии }
J_C(\pi) \leq d,
\end{equation}
где
\begin{equation}
J_C(\pi) =
\mathbb{E}_{\tau \sim \pi}
\left[
\sum_{t=0}^{T} \gamma_c^t C(s_t,a_t)
\right].
\end{equation}

Постановка CMDP является удобной, поскольку позволяет описывать компромисс между достижением цели и соблюдением ограничений в единой математической форме~\cite{altman1999constrained}. Однако на практике решение такой задачи затруднено несколькими факторами. Во-первых, агент обычно не знает функцию переходов среды и должен исследовать пространство состояний. Во-вторых, для построения безопасной стратегии агенту необходимо получать информацию о небезопасных состояниях, что само по себе может требовать посещения опасных областей. В-третьих, в задачах с редкой наградой и большим горизонтом агент может не обнаружить успешные траектории, если ограничения слишком сильно подавляют исследование.

Таким образом, центральная проблема безопасного обучения с подкреплением заключается в необходимости совместить два противоречивых требования: достаточно активное исследование среды и ограничение числа небезопасных действий. Эта проблема далее называется проблемой безопасного исследования.

\section{Методы оптимизации стратегии с ограничениями}
\label{sec:constrained_policy_optimization}

Наиболее распространённый класс методов безопасного обучения с подкреплением основан на прямой оптимизации стратегии с учётом ограничений. В таких подходах исходная задача ограниченного марковского процесса принятия решений преобразуется в задачу оптимизации, где агент должен одновременно максимизировать ожидаемую награду и удерживать ожидаемую стоимость нарушений ниже заданного порога~\cite{gu2022review,altman1999constrained}.

Одним из базовых подходов является метод множителей Лагранжа. В этом случае задача с ограничениями заменяется оптимизацией лагранжиана:
\begin{equation}
\mathcal{L}(\pi,\lambda)
=
J_R(\pi) - \lambda (J_C(\pi) - d),
\end{equation}
где $\lambda \geq 0$ --- множитель Лагранжа, определяющий вес штрафа за нарушение ограничения, $J_R(\pi)$ --- ожидаемая награда стратегии, $J_C(\pi)$ --- ожидаемая стоимость нарушений, а $d$ --- допустимый бюджет стоимости. Если ожидаемая стоимость превышает допустимый бюджет, значение $\lambda$ увеличивается, что делает стратегию более осторожной. Если ограничение выполняется с запасом, значение $\lambda$ уменьшается, позволяя агенту уделять больше внимания максимизации награды.

Лагранжевы методы получили широкое распространение, поскольку они относительно просто встраиваются в существующие алгоритмы обучения с подкреплением. К этому направлению относятся PPO-Lagrangian, TRPO-Lagrangian, TD3-Lagrangian и SAC-Lagrangian~\cite{ray2019benchmarking,stooke2020pidlag,yang2021wcsac}, а также другие варианты, использующие двойственную оптимизацию для баланса между наградой и стоимостью. Их достоинство состоит в том, что они не требуют явного знания динамики среды и могут применяться в задачах, где функция стоимости задана только через опыт взаимодействия.

К близкому направлению относится CPO~\cite{achiam2017cpo}, где обновление стратегии строится так, чтобы ограничить изменение стратегии и одновременно учитывать ограничение на стоимость. Такой подход позволяет получить более контролируемое обновление по сравнению с простыми лагранжевыми методами, однако требует более сложных вычислений, связанных с аппроксимацией доверительной области и использованием информации второго порядка, что ограничивает масштабируемость метода в больших непрерывных задачах. Дальнейшие работы стремились сделать оптимизацию с ограничениями более устойчивой и вычислительно простой. Методы CUP~\cite{cup} и FOCOPS~\cite{focops} используют специальные аппроксимации задачи с ограничениями, позволяющие получать более стабильные обновления стратегии. SAUTE RL~\cite{sootla2022saute} переносит информацию о допустимом бюджете стоимости в состояние агента, превращая задачу с ограничениями в обычную задачу обучения с подкреплением в расширенном пространстве состояний. AutoCost~\cite{he2023autocost} рассматривает автоматическое построение функции стоимости, что важно в случаях, когда корректное ручное задание стоимости затруднено.

Несмотря на различия, эти методы имеют общее ограничение: они в основном регулируют безопасность через оптимизацию ожидаемой накопленной стоимости, но не обеспечивают строгой безопасности каждого действия во время исполнения. Стратегия может нарушать ограничения на ранних этапах обучения или в состояниях, плохо представленных в обучающей выборке. Кроме того, лагранжева оптимизация может быть нестабильной: изменение множителя Лагранжа чувствительно к шуму в оценках стоимости, к изменению текущей стратегии и к масштабу награды и стоимости.

Ещё одна проблема состоит в недостаточном исследовании среды. В задачах с редкой наградой и большим горизонтом безопасное ограничение может подавлять исследование до того, как агент обнаружит успешные траектории. В навигационных задачах это часто приводит к консервативному поведению: агент избегает потенциально опасных областей, но вместе с этим не достигает цели. Поэтому методы оптимизации стратегии с ограничениями являются универсальным инструментом безопасного обучения, но не полностью решают проблему безопасного исследования в задачах с большим горизонтом, частичной наблюдаемостью и сложной структурой целей.

\section{Методы безопасного управления и фильтрации действий}
\label{sec:safety_filters}

Другой класс методов безопасного обучения использует внешние механизмы проверки и коррекции действий агента. В таких подходах обучаемая стратегия предлагает действие, после чего специальный модуль безопасности проверяет его допустимость. Если действие признаётся опасным, оно заменяется на ближайшее безопасное действие, блокируется или корректируется перед выполнением в среде. Такой подход отделяет обучение целевой стратегии от механизма обеспечения безопасности на этапе исполнения.

К этому направлению относятся слои безопасности, методы экранирования, управляющие барьерные функции и методы на основе функций Ляпунова~\cite{dalal2018safe,sac-rcbf,koller2018learning,gu2022review}. Слой безопасности обычно решает вспомогательную задачу оптимизации: найти действие, наиболее близкое к действию стратегии, но удовлетворяющее локальному ограничению безопасности. Управляющие барьерные функции и функции Ляпунова формулируют безопасность через специальные функции состояния, которые ограничивают допустимые переходы системы и позволяют получать более строгий контроль поведения.

Достоинство таких методов состоит в том, что они могут усиливать безопасность непосредственно во время исполнения стратегии. Если известны динамика системы, безопасная область и аналитическая форма ограничений, фильтр может предотвращать опасные действия независимо от ошибок обучаемой стратегии. Поэтому такие методы особенно естественны для робототехники и управления, где часто доступны физические модели, ограничения на управление или геометрическое описание опасных областей.

Однако эта же особенность ограничивает применимость методов фильтрации. Во многих задачах безопасного обучения агент не имеет точной модели динамики, полной информации о состоянии среды или аналитически заданной функции стоимости. Кроме того, методы на основе барьерных функций и слоёв безопасности часто требуют дифференцируемой функции стоимости и решения задачи оптимизации на каждом шаге управления. Это увеличивает вычислительные затраты во время исполнения и затрудняет применение таких методов в сложных, частично наблюдаемых и стохастических средах.

К близкому направлению относятся методы экранирования, в которых небезопасные действия отклоняются или заменяются во время выполнения стратегии. Во многих работах экран строится на основе формальных правил, например темпоральной логики, автоматов или заранее заданных спецификаций безопасности~\cite{hunt2021verifiably,yangsafe}. Такой подход позволяет явно контролировать поведение агента, но требует ручного задания ограничений и часто плохо масштабируется на непрерывные пространства состояний и действий.

Более новые методы пытаются уменьшить зависимость от ручных правил. Например, ADVICE~\cite{advice} рассматривает экранирование в более общей постановке: вместо аналитической модели динамики или формального описания опасных состояний используется обучаемый дискриминатор безопасности, построенный по траекториям. Такой подход снижает требования к априорному знанию среды, но всё равно сохраняет отдельный модуль фильтрации во время исполнения, что увеличивает сложность системы и вычислительные затраты.

В целом методы безопасного управления и фильтрации действий обеспечивают более прямой контроль исполнения, чем чистая оптимизация стратегии с ограничениями. Однако они обычно требуют дополнительной информации о среде: модели динамики, дифференцируемой функции стоимости, полной наблюдаемости, формальных правил или возможности решать задачу оптимизации на каждом шаге. Поэтому их применение ограничено в задачах, где агент действует в условиях частичной наблюдаемости, неизвестной динамики и недифференцируемых или неявно заданных ограничений безопасности.

Для задач, рассматриваемых далее в диссертации, важны именно эти ограничения существующих подходов. Методы оптимизации стратегии с ограничениями являются достаточно общими, но страдают от нестабильности и слабого исследования в задачах с большим горизонтом. Методы фильтрации и безопасного управления могут давать более строгий контроль действий, но часто требуют сильных предположений о среде. Это мотивирует рассмотрение подходов, которые сохраняют общность оптимизационных методов, но улучшают исследование среды и оценку риска без необходимости использовать отдельный фильтр безопасности во время исполнения.

\section{Методы безопасного обучения на основе моделей мира}
\label{subsec:safe_mbrl}

Несмотря на развитие методов обучения с подкреплением без модели мира, их применение в задачах с жёсткими ограничениями остаётся затруднённым. Основная причина состоит в том, что такие методы требуют большого числа непосредственных взаимодействий со средой. В безопасном обучении это особенно критично: на ранних этапах агент действует методом проб и ошибок, поэтому увеличение числа обучающих эпизодов обычно сопровождается ростом числа нарушений ограничений безопасности. В результате даже алгоритмы, учитывающие ограничения в процессе оптимизации стратегии, могут накапливать значительные штрафы во время обучения.

Одним из способов уменьшить число небезопасных взаимодействий является использование моделей мира. Методы обучения с подкреплением на основе модели мира строят приближённую модель динамики среды, которая предсказывает последствия действий агента. Такая модель может использоваться для генерации синтетического опыта, обновления стратегии по воображаемым траекториям или предварительной оценки действий перед их выполнением в реальной среде~\cite{polydoros2017survey,moerland2023model}. В задачах безопасного обучения модель мира особенно важна, поскольку позволяет заранее оценивать не только ожидаемую награду, но и будущую стоимость кандидатных действий.

С точки зрения использования модели мира можно выделить два основных направления. В первом случае модель применяется непосредственно при выборе действия: агент рассматривает несколько кандидатных последовательностей действий, оценивает их в модели мира и выполняет первое действие из выбранной последовательности. К этому направлению относятся методы на основе управления с прогнозирующей моделью и планирования с ограничениями~\cite{camacho2007model,koller2018learning,wabersich2021predictive,liu2020mbrce,wen2018cce}. Во втором случае модель используется для генерации синтетических траекторий и обновления стратегии без обязательного планирования на каждом шаге~\cite{hafner2019dreamer,jayant2022model,thomas2021safe}. Оба направления повышают эффективность использования данных, но имеют разные ограничения: планирование при выборе действия может быть вычислительно дорогим и учитывать только короткий горизонт, а обучение на воображаемых траекториях чувствительно к ошибкам модели.

Одним из ранних направлений безопасного обучения с подкреплением на основе модели мира является безопасное планирование с использованием ограничений. Метод constrained cross-entropy~\cite{wen2018cce} расширяет метод кросс-энтропии для задач с ограничениями: кандидатные траектории оцениваются по награде и стоимости, а в элитное множество попадают только траектории, удовлетворяющие ограничению на стоимость. Похожую идею развивает MBRCE~\cite{liu2020mbrce}, где используются ансамбли моделей и устойчивый поиск действий через метод кросс-энтропии. Такие методы позволяют отбрасывать потенциально опасные последовательности действий до их выполнения, однако их качество зависит от точности предсказаний модели и от выбранного горизонта планирования.

Другая группа работ использует модель мира для обновления стратегии. MBPPO-Lag~\cite{jayant2022model} расширяет PPO-Lagrangian синтетическими траекториями, полученными из ансамбля моделей мира. LAMBDA~\cite{as2022lambda} сочетает лагранжеву оптимизацию с байесовской моделью мира и развивает идеи Dreamer для безопасного обучения. Safe SLAC~\cite{hogewind2022safe} добавляет лагранжевый механизм в латентную стохастическую модель SLAC. Эти методы показывают, что модели мира могут повышать эффективность безопасного обучения, однако часто не используют модель мира для непосредственного безопасного планирования действий.

Существенное влияние на развитие обучения с подкреплением на основе модели оказали методы Dreamer и TD-MPC. Dreamer обучает скрытую модель динамики и использует воображаемые траектории для оптимизации поведения~\cite{hafner2019dreamer,hafner2023dreamerv3}. TD-MPC и TD-MPC2 используют неявные модели мира и совмещают планирование в скрытом пространстве с функциями оценки полезности~\cite{hansen2022tdmpc,hansen2024tdmpc2}. Эти подходы показали высокую эффективность в задачах непрерывного управления, но в исходном виде не предназначены для соблюдения ограничений безопасности.

К этому классу методов относится SafeDreamer~\cite{safedreamer}, который развивает идеи DreamerV3 и объединяет лагранжевую оптимизацию, планирование при выборе действия и обучение по воображаемым траекториям. В отличие от ряда предыдущих методов, SafeDreamer поддерживает как низкоразмерные, так и визуальные наблюдения, использует модель мира для безопасного планирования и обучает стратегию с учётом ограничений. Такой подход показывает сильные результаты в задачах SafetyGymnasium~\cite{ji2023safetygymnasium}. Однако архитектуры семейства Dreamer могут быть вычислительно тяжёлыми, а планирование в модели мира остаётся ограниченным коротким горизонтом.

Для безопасного обучения с подкреплением на основе модели мира критически важна оценка будущей стоимости. Если модель недооценивает риск, агент может выбирать небезопасные траектории, поскольку опасные последствия не всегда проявляются в пределах короткого горизонта планирования. Если модель переоценивает риск, агент становится чрезмерно осторожным и прекращает активное исследование среды. Поэтому важно не только предсказывать мгновенную стоимость, но и строить функцию долгосрочной стоимости, оценивающую ожидаемые будущие нарушения.

Эта проблема связана с ограничениями классического планирования на конечном горизонте. При коротком горизонте оптимизация может учитывать только ближайшие последствия действий, если после его окончания не используется терминальная оценка полезности или стоимости. В задачах оптимизации награды эта проблема частично решается добавлением функции оценки полезности в конце горизонта~\cite{sikchi2022loop,moerland2023model}. Однако для безопасного обучения аналогичная проблема возникает и для стоимости: необходимо учитывать не только будущую награду, но и долгосрочный риск нарушений безопасности.

Отдельную роль играет неопределённость модели. Модель мира обучается по конечному набору данных и может ошибаться в областях состояния, которые агент редко посещал. Поэтому в безопасном обучении с подкреплением на основе модели мира часто используются ансамбли моделей или функций стоимости~\cite{liu2020mbrce,jayant2022model,as2022lambda}. Ансамблевые оценки позволяют уменьшить влияние ошибок отдельной модели, однако выбор способа агрегации становится самостоятельной проблемой. Оптимистичная агрегация может приводить к недооценке риска, а чрезмерно пессимистичная~--- к избыточно осторожному поведению.

\section{Иерархическое обучение с подкреплением и задачи с большим горизонтом}
\label{sec:hrl}

Иерархическое обучение с подкреплением рассматривает задачу принятия решений на нескольких уровнях абстракции. Основная идея состоит в том, что агент выбирает не только примитивные действия, но и более продолжительные элементы поведения: подзадачи, навыки или промежуточные подцели. Такая декомпозиция позволяет заменить длинную последовательность примитивных действий более короткой последовательностью высокоуровневых решений и тем самым уменьшить эффективный горизонт задачи~\cite{pateria2021hrl}.

В общем виде иерархический агент включает верхнеуровневую стратегию и одну или несколько низкоуровневых стратегий. Верхнеуровневая стратегия выбирает абстрактное действие: подзадачу, навык или подцель. Низкоуровневая стратегия исполняет это решение в течение нескольких шагов и выдаёт примитивные действия в среду. За счёт этого возникает временная абстракция: одно решение верхнего уровня соответствует нескольким действиям нижнего уровня. Такая структура особенно полезна в задачах с редкой наградой, где одной стратегии трудно случайно обнаружить успешную последовательность действий.

В литературе можно выделить два распространённых способа задания высокоуровневых действий. Первый способ основан на навыках или опциях. Опция задаётся стратегией выполнения, условием начала и условием завершения. В этом случае верхнеуровневая стратегия выбирает, какой навык следует активировать, а выбранный навык некоторое время управляет агентом. К этому направлению относятся классические работы по опциям~\cite{sutton1999options}, а также методы, обучающие опции совместно с основной стратегией, например Option-Critic~\cite{bacon2017optioncritic}. Преимущество такого подхода состоит в том, что навыки могут представлять переиспользуемые фрагменты поведения. Однако в задачах с редкой наградой обучение полезных навыков с нуля может быть затруднено, если отсутствует дополнительный сигнал о том, какие навыки должны быть сформированы.

Второй способ основан на подцелях. В этом случае верхнеуровневая стратегия выбирает состояние, признак состояния или вектор в некотором пространстве признаков, а низкоуровневая стратегия обучается достигать выбранную подцель. Методы HAC~\cite{levy2017hierarchical}, HiRO~\cite{nachum2018data} и HRAC~\cite{hrac} используют такую структуру для разбиения исходной задачи на более короткие подзадачи. Верхний уровень работает на более грубой временной шкале и выбирает подцели, а нижний уровень действует чаще и решает локальную задачу достижения заданной подцели.

Одной из ключевых проблем иерархического обучения является нестационарность. Пока низкоуровневая стратегия обучается, её способность достигать подцелей меняется. Поэтому один и тот же выбор верхнего уровня может приводить к разным фактическим переходам в разные моменты обучения. Метод HiRO уменьшает эту проблему с помощью переобозначения подцелей в собранных траекториях, согласуя высокоуровневое обучение с фактически достигнутыми состояниями~\cite{nachum2018data}. Метод HAC также использует переобозначение подцелей и позволяет строить многоуровневую иерархию~\cite{levy2017hierarchical}.

Другой важной проблемой является достижимость подцелей. Если верхнеуровневая стратегия выбирает подцели, которые нижний уровень не может достичь, обучение становится нестабильным, а собранный опыт плохо соответствует задаче верхнего уровня. В HRAC для ограничения пространства подцелей используется информация о смежности состояний: подцели выбираются так, чтобы они были достижимы низкоуровневой стратегией за ограниченное число шагов~\cite{hrac}. Тем самым уменьшается вероятность выбора недостижимых подцелей и повышается устойчивость обучения.

Несмотря на эффективность иерархической декомпозиции в задачах с большим горизонтом, классические методы иерархического обучения обычно не учитывают требования безопасности. Они предполагают, что достижимая подцель является допустимой, либо полагаются на тщательно заданную функцию награды. В безопасном обучении это предположение недостаточно: подцель может быть достижимой с точки зрения динамики, но её достижение может требовать прохождения через опасные области или накопления слишком большой стоимости. Поэтому достижимость и безопасность подцелей являются разными свойствами.

Интеграция безопасности в иерархическое обучение остаётся менее изученной задачей. Один вариант состоит в том, чтобы использовать безопасную коррекцию действий на нижнем уровне. Например, методы, объединяющие Safety Layer и HiRO~\cite{safetylayerhiro}, корректируют действия низкоуровневой стратегии с помощью дополнительного модуля безопасности. Такой подход может уменьшать нарушения, однако требует дифференцируемой функции стоимости и дополнительной оптимизации во время исполнения. Кроме того, он сохраняет зависимость от того, насколько точно низкоуровневая стратегия достигает заданных подцелей.

Другой вариант состоит в использовании планирования или функций Ляпунова для построения безопасных траекторий между подцелями. Например, LyapunovRRT~\cite{lyapunovrrt} объединяет планирование на графе с нейронными функциями Ляпунова для генерации безопасных путей. Такие методы могут обеспечивать более строгий контроль безопасности, но требуют полной информации о состоянии среды, положении препятствий или области безопасных состояний. Это ограничивает их применимость в частично наблюдаемых задачах и в средах, где опасные области заранее неизвестны.

Существуют также методы безопасного иерархического обучения, основанные на заранее заданных опциях или предметно-специфических навыках, например IAHRL~\cite{iahrl}. Такие подходы используют иерархию для ускорения обучения и уменьшения горизонта задачи, но часто требуют ручного задания набора навыков или структуры подзадач. Поэтому они хуже подходят для общей постановки, где агент должен обучаться без заранее заданной иерархии поведения.

Таким образом, иерархическое обучение с подкреплением эффективно решает проблему большого горизонта за счёт временной абстракции, подцелей и навыков. Однако само наличие иерархии не обеспечивает безопасность. Для безопасных задач необходимо учитывать не только достижимость подцелей, но и стоимость их достижения. Существующие безопасные иерархические методы часто требуют полной наблюдаемости, дифференцируемых функций стоимости, заранее заданных навыков или дополнительной фильтрации действий во время исполнения. Это мотивирует разработку подходов, в которых иерархическая структура используется для улучшения исследования, а ограничения безопасности учитываются в процессе обучения без необходимости полагаться на ручные правила или отдельный фильтр действий на этапе исполнения.

\section{Безопасное иерархическое обучение с подкреплением}
\label{sec:safe_hrl}

Интеграция иерархического обучения с подкреплением и ограничений безопасности остаётся менее изученным направлением по сравнению с обычным безопасным обучением и классическим иерархическим обучением. Эта область важна, поскольку многие робототехнические задачи одновременно обладают большим горизонтом, редкой наградой и требованиями к безопасности. В таких задачах агенту необходимо активно исследовать среду для обнаружения успешных траекторий, однако ограничения безопасности могут подавлять именно те действия, которые нужны для исследования.

Классическая иерархическая декомпозиция уменьшает эффективный горизонт задачи за счёт промежуточных подцелей или навыков. Однако сама по себе такая декомпозиция не гарантирует безопасность. Подцель может быть достижимой с точки зрения динамики, но её достижение может требовать прохождения через опасные области или накопления слишком большой стоимости. Поэтому в безопасном иерархическом обучении необходимо учитывать не только достижимость подцелей, но и стоимость их достижения.

В литературе можно выделить три близких метода безопасного иерархического обучения с подкреплением: SafetyLayer + HiRO~\cite{safetylayerhiro}, LyapunovRRT~\cite{lyapunovrrt} и IAHRL~\cite{iahrl}. Эти методы используют иерархическую структуру для решения задач с большим горизонтом, но различаются способом обеспечения безопасности и требованиями к информации о среде.

Метод SafetyLayer + HiRO~\cite{safetylayerhiro} объединяет иерархическую стратегию HiRO с дополнительным слоем безопасности. Верхнеуровневая стратегия задаёт подцели, а низкоуровневая стратегия выбирает действия для их достижения. Перед выполнением действие проверяется модулем безопасности и при необходимости корректируется. Такой подход может уменьшать число нарушений безопасности, однако требует дифференцируемой функции стоимости и вводит дополнительные вычисления на этапе исполнения. Кроме того, безопасность обеспечивается в основном через коррекцию примитивных действий, тогда как безопасность самих подцелей остаётся отдельной проблемой.

Метод LyapunovRRT~\cite{lyapunovrrt} сочетает планирование RRT* с контроллером на основе функций Ляпунова. Он строит безопасные траектории между состояниями и может использоваться для генерации безопасных путей в иерархической постановке. Однако такой подход требует полной информации о состоянии среды, положении препятствий и динамике системы для вычисления областей притяжения. Это ограничивает его применимость в частично наблюдаемых задачах и в средах, где агент не имеет доступа к полной карте или точной модели динамики.

Метод IAHRL~\cite{iahrl} использует иерархическую структуру с заранее определёнными опциями. Наличие таких опций позволяет ускорить обучение и улучшить исследование в конкретных доменах. Однако метод зависит от ручного проектирования структуры поведения и плохо переносится на новые среды, где набор полезных навыков заранее неизвестен.

Сравнение рассмотренных методов и метода SG-Safe приведено в таблице~\ref{tab:safe_hrl_comparison}. В таблицу включены как универсальные методы безопасного обучения с подкреплением, так и специализированные безопасные иерархические методы. Такое сравнение показывает, какие дополнительные предположения требуются разным подходам: полная наблюдаемость, явная или дифференцируемая функция стоимости, дополнительный контроль безопасности на этапе исполнения или ручная разработка навыков.

\begin{table}[ht]
\centering
\caption{Сравнение методов безопасного обучения с подкреплением и безопасного иерархического обучения. FS~--- не требуется полное состояние среды; EC~--- не требуется явная функция стоимости; DC~--- не требуется дифференцируемая функция стоимости; RC~--- не требуется дополнительный контроль безопасности на этапе исполнения; LH~--- метод специально ориентирован на задачи с большим горизонтом; NSD~--- не требуется ручная разработка навыков или опций.}
\label{tab:safe_hrl_comparison}
\resizebox{0.68\linewidth}{!}{
\begin{tabular}{@{}l*{6}{c}@{}}
\toprule
Method & FS & EC & DC & RC & LH & NSD \\
\midrule
Lagrangian Safe RL~\cite{yang2021wcsac,stooke2020pidlag,ray2019benchmarking}
& \checkmark & \xmark & \checkmark & \checkmark & \xmark & \checkmark \\
CPO~\cite{achiam2017cpo}
& \checkmark & \xmark & \checkmark & \checkmark & \xmark & \checkmark \\
CUP~\cite{cup}
& \checkmark & \xmark & \checkmark & \checkmark & \xmark & \checkmark \\
FOCOPS~\cite{focops}
& \checkmark & \xmark & \checkmark & \checkmark & \xmark & \checkmark \\
SafeDreamer~\cite{safedreamer}
& \checkmark & \xmark & \checkmark & \xmark & \xmark & \checkmark \\
Safety Layer~\cite{dalal2018safe}
& \xmark & \xmark & \xmark & \xmark & \xmark & \checkmark \\
CBF-based Safe RL~\cite{sac-rcbf}
& \xmark & \xmark & \xmark & \xmark & \xmark & \checkmark \\
Shielding~\cite{hunt2021verifiably,yangsafe}
& \xmark & \checkmark & \checkmark & \xmark & \xmark & \checkmark \\
ADVICE~\cite{advice}
& \checkmark & \checkmark & \checkmark & \xmark & \xmark & \checkmark \\
SafetyLayer + HiRO~\cite{safetylayerhiro}
& \xmark & \xmark & \xmark & \xmark & \checkmark & \checkmark \\
LyapunovRRT~\cite{lyapunovrrt}
& \xmark & \xmark & \xmark & \xmark & \checkmark & \checkmark \\
IAHRL~\cite{iahrl}
& \checkmark & \checkmark & \checkmark & \checkmark & \checkmark & \xmark \\
SG-Safe
& \checkmark & \checkmark & \checkmark & \checkmark & \checkmark & \checkmark \\
\bottomrule
\end{tabular}
}
\end{table}

Таким образом, универсальные методы безопасного обучения с подкреплением, такие как Lagrangian Safe RL, CPO, CUP и FOCOPS, не требуют полной информации о состоянии среды и не используют дополнительный контроль безопасности на этапе исполнения, однако они не содержат специального механизма для решения задач с большим горизонтом. Методы безопасного управления и фильтрации действий, такие как Safety Layer, CBF-based Safe RL и Shielding, усиливают безопасность во время исполнения, но требуют дополнительных предположений о среде или отдельного модуля контроля действий. Метод ADVICE ослабляет часть этих требований, но также использует фильтрацию действий на этапе исполнения.

Существующие безопасные иерархические методы ближе к рассматриваемой постановке, но также имеют ограничения. SafetyLayer + HiRO требует дифференцируемой функции стоимости и дополнительной коррекции действий во время исполнения. LyapunovRRT требует полной наблюдаемости и информации о динамике среды. IAHRL ориентирован на задачи с большим горизонтом и не требует дополнительного фильтра действий, однако зависит от заранее заданных опций, что ограничивает переносимость метода на новые задачи. На этом фоне SG-Safe направлен на использование иерархической структуры для улучшения исследования в задачах с большим горизонтом без полной наблюдаемости, ручной разработки навыков и дополнительного контроля безопасности на этапе исполнения.

\section{Мультимодальные агенты и естественно-языковые инструкции}
\label{sec:multimodal_agents}

Современные автономные агенты всё чаще рассматриваются не только как системы управления, но и как системы, взаимодействующие с человеком через естественный язык. В таких задачах агент должен понимать текстовую инструкцию, сопоставлять её с объектами и состояниями среды, а затем выбирать действия, приводящие к выполнению инструкции. Это направление связано с language grounding, embodied AI и мультимодальным обучением с подкреплением.

\subsection{Бенчмарки безопасного обучения с подкреплением}
\label{subsec:safe_rl_benchmarks}

Развитие безопасного обучения с подкреплением во многом определяется наличием бенчмарков, позволяющих сравнивать методы в воспроизводимых условиях. В прикладных задачах обучение с подкреплением используется в управлении транспортом, робототехнике, энергетике, автономных системах и других областях, где ошибки агента могут приводить к существенным негативным последствиям~\cite{campbell2010autonomous,liang2019traffic,wei2017energy,kayhan2023scheduling}. Поэтому для Safe RL важны среды, в которых можно одновременно оценивать достижение цели и соблюдение ограничений безопасности.

Ранние среды для исследования безопасности часто строились на простых дискретных задачах. Примером являются AI Safety Gridworlds~\cite{leike2017safety_gridworlds}, где рассматриваются задачи безопасного прерывания, устойчивости к сдвигу распределения и предотвращения нежелательного поведения. Такие среды удобны для анализа отдельных свойств алгоритмов, однако они имеют упрощённую динамику, дискретное пространство состояний и действий, а также слабо отражают сложности реальных задач управления.

Другой класс сред связан с адаптацией существующих RL-бенчмарков под безопасное обучение. В частности, в MiniGrid~\cite{chevalier2024minigrid} и его модификациях опасные клетки или области используются как источники штрафов, которых агент должен избегать~\cite{wachi2021safe,wachi2024safe_exploration}. Такие среды позволяют проверять базовые способности к безопасной навигации, но остаются детерминированными, маломерными и в основном сводятся к поиску пути в двумерной сетке.

Для задач непрерывного управления широко используется Safety Gym~\cite{ray2019benchmarking}. В нём агенты различной морфологии, такие как Point, Car, Ant и Doggo, решают навигационные задачи при наличии опасных зон, ваз, кнопок и других объектов. Safety Gym стал одним из основных бенчмарков для сравнения методов безопасного исследования и лагранжевых подходов к учёту ограничений. Позднее Safety Gymnasium~\cite{ji2023safetygymnasium} расширил этот набор задач, добавив больше агентов, сценариев и вариантов постановки, включая задачи для многоагентного безопасного обучения.

Несмотря на широкое использование Safety Gym и Safety Gymnasium, эти бенчмарки имеют ряд ограничений. Во-первых, многие задачи остаются преимущественно навигационными: агент должен двигаться к цели и избегать опасных областей. Во-вторых, опасные объекты часто статичны или движутся по заранее заданным траекториям. В-третьих, ограничения безопасности задаются фиксированной численной функцией стоимости, а не гибким описанием поведения, которое может изменяться от задачи к задаче. Поэтому такие среды хорошо подходят для проверки классического Safe RL, но хуже отражают сценарии, где требования безопасности задаются человеком в более сложной форме.

Отдельное направление связано с безопасным управлением в робототехнических и физических симуляторах. Safe-Control-Gym~\cite{yuan2022safecontrolgym} включает задачи управления cartpole и quadrotor с возмущениями динамики и ориентирован на проверку методов безопасного управления и sim-to-real переноса. RWRL~\cite{dulac2020realworldrl} добавляет ограничения и возмущения в действия, наблюдения и физические параметры для робототехнических задач. Safe MAMuJoCo, Safe MARobosuite и Safe MAIG~\cite{gu2023safe_multiagent} предназначены для безопасного многоагентного обучения и робототехнического управления. Эти среды ближе к физическому управлению, но часто требуют дорогого моделирования и в основном используют проприоцептивные признаки, а не сложное визуальное восприятие.

Для автономного вождения и навигации применяются симуляторы CARLA~\cite{dosovitskiy2017carla}, TORCS~\cite{wymann2000torcs} и MetaDrive~\cite{li2022metadrive}. В таких средах безопасность обычно учитывается через штрафы за столкновения, выезд за пределы полосы или другие нежелательные события~\cite{wang2018autonomous,nehme2023safe,hossain2023carla}. Эти бенчмарки позволяют моделировать более реалистичные сценарии движения, однако они ориентированы на конкретные прикладные домены и не всегда удобны для систематического анализа общих алгоритмов Safe RL.

Более современным примером является HASARD~\cite{hasard}, построенный на основе ViZDoom~\cite{kempka2016vizdoom}. В отличие от многих предыдущих бенчмарков, HASARD ориентирован на безопасное обучение с эгоцентрическим визуальным восприятием в трёхмерной стохастической среде. Он включает несколько уровней сложности, частичную наблюдаемость и задачи, требующие не только навигации, но и позиционирования, распознавания целей и предсказания движения объектов. При этом использование ViZDoom и Sample-Factory~\cite{petrenko2020samplefactory} позволяет проводить обучение быстрее, чем в тяжёлых физически точных симуляторах. Однако HASARD, как и большинство Safe RL бенчмарков, не проверяет способность агента понимать естественно-языковые инструкции и ограничения.

Таким образом, существующие Safe RL бенчмарки можно условно разделить на несколько групп: простые двумерные среды для анализа отдельных аспектов безопасности, среды непрерывного управления типа Safety Gym, робототехнические и физические симуляторы, а также визуальные 3D-среды. Несмотря на различия, большинство из них задаёт безопасность через фиксированную функцию стоимости. Такая постановка удобна для классического Safe RL, но ограничивает переход к задачам, где человек задаёт цель и ограничения в форме естественного языка.

\subsection{Бенчмарки с естественно-языковыми ограничениями}
\label{subsec:language_safe_rl_benchmarks}

Переход от фиксированных численных ограничений к естественно-языковым инструкциям существенно усложняет задачу безопасного обучения с подкреплением. В классической целеобусловленной постановке цель задаётся как элемент пространства целей ($g \in \mathcal{G}$), а стратегия имеет вид
\begin{equation}
a_t \sim \pi(a_t \mid o_t, g),
\end{equation}
где $o_t$~--- наблюдение агента. В мультимодальной постановке цель задаётся текстовой инструкцией $I_g$, поэтому стратегия принимает вид
\begin{equation}
a_t \sim \pi(a_t \mid o_t, I_g).
\end{equation}
В этом случае агент должен сопоставить текст с объектами, действиями и состояниями среды.

Для безопасного обучения с естественно-языковыми ограничениями вводится дополнительная инструкция $I_c$, задающая правило безопасности. Тогда полная инструкция имеет вид $I=(I_g,I_c)$, где $I_g$ задаёт целевую задачу, а $I_c$ задаёт ограничение. Стратегия должна учитывать оба компонента:
\begin{equation}
a_t \sim \pi(a_t \mid o_t, I_g, I_c).
\end{equation}

Такая постановка отличается как от классического Safe RL, так и от стандартного instruction following. В обычном Safe RL ограничения, как правило, задаются в математической или логической форме, например через функцию стоимости или логическое правило. Это требует экспертного задания ограничений и затрудняет использование таких систем конечными пользователями. Естественный язык, напротив, является более доступным способом задания ограничений: человек может сформулировать правило вида ``не заходи в опасную область до выполнения другого действия'' без явного задания соответствующей формулы стоимости.

При этом естественно-языковые ограничения не совпадают с обычными инструкциями на выполнение задачи. В задачах instruction following инструкция обычно задаёт, что агент должен сделать, и тем самым непосредственно определяет целевое поведение~\cite{macmahon2006walk,chen2011learning,artzi2013weakly,misra2017mapping,tellex2011understanding,andreas2015alignment}. В безопасном обучении с языковыми ограничениями текстовое правило, напротив, задаёт, чего агент не должен делать, оставаясь ортогональным функции вознаграждения. Агент должен самостоятельно искать поведение, максимизирующее награду, но при этом соблюдать заданное текстовое ограничение.

Смежное направление связано с vision-and-language navigation и embodied AI, где агент должен выполнять навигационные инструкции в визуальной или трёхмерной среде~\cite{anderson2018vision,misra2018mapping,hermann2020directions,hao2020pretraining,chen2019touchdown,de_vries2018talk}. Эти работы важны для grounding естественного языка в наблюдениях среды, однако в большинстве таких постановок язык используется для задания цели, а не независимого ограничения безопасности. Поэтому они не проверяют способность агента отделять достижение цели от соблюдения запретов.

Одной из первых работ, специально рассматривающих безопасное обучение с естественно-языковыми ограничениями, является HAZARDWORLD~\cite{hazardworld}. Этот бенчмарк включает две группы сред: HAZARD WORLD grid на основе MiniGrid~\cite{minigrid} и HAZARD WORLD robot на основе Safety Gym~\cite{ray2019benchmarking}. В HAZARDWORLD вводится разделение на этапы safety training и safety evaluation: во время обучения агент должен научиться интерпретировать текстовые ограничения, а во время оценки~--- переносить это понимание на новые задачи с другими функциями вознаграждения.

В HAZARDWORLD рассматриваются три основных типа естественно-языковых ограничений. Бюджетные ограничения ограничивают число посещений опасных состояний, например запрещают заходить в лаву более заданного числа раз. Отношенческие ограничения определяют опасные состояния через пространственные отношения с объектами среды, например требуют держаться на определённом расстоянии от опасного объекта. Последовательные ограничения делают состояние опасным в зависимости от истории взаимодействия, например запрещают наступать на воду после посещения травы. Такие ограничения требуют от агента не только распознавания текущего состояния, но и учёта истории, порядка действий и отношений между объектами.

Для демонстрации этой постановки в HAZARDWORLD предложен метод POLCO, в котором обучение представлений для текстовых ограничений отделено от обучения стратегии. Сначала constraint interpreter преобразует языковое ограничение в представление запрещённых состояний, после чего strategy network использует это представление вместе с наблюдением для выбора действия. Такая факторизация позволяет переносить понимание ограничений на новые задачи, где функция вознаграждения отличается от использованной при обучении.

Несмотря на значимость HAZARDWORLD, этот бенчмарк сохраняет ограничения исходных сред. MiniGrid-версия имеет простую дискретную динамику и маломерные наблюдения. SafetyGym-версия ближе к непрерывному управлению, но остаётся в основном навигационной задачей. Кроме того, естественный язык используется преимущественно для задания ограничений, тогда как конечная цель часто задаётся координатами, фиксированной структурой среды или функцией вознаграждения. Это ограничивает применимость такой постановки к сценариям, где человек формулирует и цель, и правило безопасности на естественном языке.

Отдельно следует отметить работы, использующие естественный язык для задания логических или структурированных ограничений. Например, языковые команды могут переводиться в логические формулы или планировочные ограничения для робототехники~\cite{howard2014natural,wang2020ltl}. Другие подходы обучают отдельные модели для проверки того, нарушает ли траектория заданное текстовое ограничение~\cite{prakash2020guiding}. Однако такие подходы часто либо сводят ограничение к дополнительному отрицательному вознаграждению, либо не рассматривают полноценную задачу безопасного обучения, где ограничение должно быть отделено от целевой функции.

Бенчмарк Caged CrafText, рассматриваемый далее в диссертации, направлен на преодоление этих ограничений. В отличие от HAZARDWORLD, он использует более сложную, стохастическую и недетерминированную среду Crafter, в которой агент должен добывать ресурсы, создавать предметы и планировать длинные последовательности действий. При этом как цель, так и ограничение задаются на естественном языке. Это требует от агента не только безопасного поведения, но и понимания инструкции, долгосрочного планирования и рассуждения о правилах, зависящих от контекста и истории взаимодействия со средой.

Таким образом, развитие бенчмарков показывает переход от классического Safe RL с фиксированными численными функциями стоимости к постановкам, где ограничения задаются естественным языком. Такой переход важен для автономных агентов, взаимодействующих с человеком: в реальных сценариях пользователь чаще задаёт правила в текстовой форме, а не в виде математической функции стоимости. Поэтому бенчмарки с естественно-языковыми целями и ограничениями необходимы для оценки агентов, которые должны одновременно понимать инструкции, планировать действия и соблюдать явно заданные правила безопасности.

\section{Выводы по главе}
\label{sec:chapter1_conclusion}

В данной главе рассмотрены основные постановки и методы безопасного обучения с подкреплением для автономных интеллектуальных агентов. В качестве базовой формальной модели использована постановка марковского процесса принятия решений с ограничениями, в которой агент должен максимизировать ожидаемую награду при соблюдении ограничения на ожидаемую стоимость нарушений безопасности.

Проведённый обзор показывает, что центральной проблемой безопасного обучения является противоречие между исследованием среды и соблюдением ограничений. С одной стороны, агенту необходимо активно исследовать среду для обнаружения успешных траекторий. С другой стороны, такое исследование может приводить к посещению опасных состояний и накоплению штрафов. Эта проблема особенно выражена в задачах с редкой наградой, большим горизонтом и частичной наблюдаемостью.

Методы оптимизации стратегии с ограничениями, включая лагранжевые подходы и методы ограниченного обновления стратегии, являются универсальными и не требуют явного знания динамики среды. Однако они обычно регулируют безопасность через ожидаемую накопленную стоимость и не обеспечивают строгой безопасности каждого действия. Кроме того, такие методы могут страдать от нестабильности двойственной оптимизации и консервативного поведения, при котором агент избегает исследования и не достигает цели.

Методы безопасного управления и фильтрации действий позволяют усиливать безопасность непосредственно на этапе исполнения. К ним относятся слои безопасности, методы экранирования, управляющие барьерные функции и подходы на основе функций Ляпунова. Их преимуществом является более прямой контроль действий агента, однако они часто требуют полной информации о состоянии среды, известной динамики, дифференцируемой функции стоимости, формальных правил или дополнительной оптимизации на каждом шаге управления. Это ограничивает их применимость в сложных частично наблюдаемых средах.

Методы безопасного обучения на основе моделей мира позволяют заранее оценивать последствия действий агента и тем самым снижать число небезопасных взаимодействий со средой. Однако их эффективность ограничивается ошибками модели, вычислительной стоимостью планирования и коротким горизонтом предсказания. Поэтому для таких методов важной проблемой становится оценка не только мгновенной стоимости, но и долгосрочного риска будущих нарушений безопасности.

Иерархическое обучение с подкреплением является естественным способом работы с задачами большого горизонта, поскольку использует временную абстракцию, подцели и навыки. Однако классические иерархические методы не учитывают безопасность подцелей: достижимая подцель может требовать небезопасной траектории. Существующие методы безопасного иерархического обучения частично решают эту проблему, но часто требуют полной наблюдаемости, дифференцируемой функции стоимости, заранее заданных навыков или дополнительного контроля действий при исполнении.

Отдельной проблемой является переход от фиксированных численных функций стоимости к естественно-языковым целям и ограничениям. В реальных сценариях пользователь может задавать не только цель, но и правило безопасности в текстовой форме. Поэтому безопасный агент должен не только оптимизировать награду и стоимость, но и связывать естественно-языковые инструкции с объектами, действиями, историей и состояниями среды.

Таким образом, проведённый обзор выделяет несколько ключевых нерешённых проблем: безопасное исследование в задачах с большим горизонтом, нестабильность оптимизации ограничений, зависимость методов фильтрации от сильных предположений о среде, необходимость оценки долгосрочной стоимости в моделях мира, безопасность подцелей в иерархическом обучении и переход к естественно-языковым ограничениям. Эти проблемы определяют направления дальнейших глав диссертации: использование моделей мира для оценки долгосрочного риска, иерархическое безопасное исследование, контроль безопасности подцелей и обучение агентов с естественно-языковыми целями и ограничениями.
